% Part: proof-theory
% Chapter: proof-search
% Section: completeness

\documentclass[../../../include/open-logic-section]{subfiles}

\begin{document}

\olfileid{pt}{ps}{cpl}

\olsection{في الاكتمال}

ونثبت الآن ضمان خوارزمية البحث عن البرهان: فإذا أجهضت أو مضت بغير توقف،
لم يوجد~!!{proof} للمتتالية الابتدائية~$\OLId{greek-variable}{Gamma} \Sequent \OLId{greek-variable}{Delta}$. وسبيل
الإثبات أن ننشئ~!!a{structure}~$\Struct M$ يتحقق فيه~$\Sat{\OLId{latin-looped}{M}}{!A}$ لكل
$!A \in \OLId{greek-variable}{Gamma}$، و$\Sat/{\OLId{latin-looped}{M}}{!A}$ لكل~$!A \in \OLId{greek-variable}{Delta}$. فتكون جميع
!!{formula}s~$\OLId{greek-variable}{Gamma}$ صادقة، وجميع~!!{formula}s~$\OLId{greek-variable}{Delta}$ كاذبة، في
$\Struct M$؛ أي تكون~$\OLId{greek-variable}{Gamma} \Sequent \OLId{greek-variable}{Delta}$ غير صالحة. ولما كان
\Log{G3c} سليمًا، امتنع لها~!!{proof}~$\OLId{greek-variable}{Gamma} \Sequent \OLId{greek-variable}{Delta}$.

ونبني~$\Struct{M}$ من زوج~$\Theta,\Xi$ من~!!{formula}s ومن مجموعة~$\OLId{latin-looped}{C}$
من~!!{constant}s. وهذه~$\Theta,\Xi$ و$\OLId{latin-looped}{C}$ نستخرجها من \emph{فرع فشل}
لتشغيل لم ينجح فيه البحث عن البرهان، إما لإجهاضه وإما لعدم توقفه. وفرع
الفشل متتالية
منتهية أو لامتناهية من المتتاليات~$\Pi_\OLId{latin-ordinary}{n} \Sequent \Lambda_\OLId{latin-ordinary}{n}$
ومجموعات~!!{constant}s~$\OLId{latin-looped}{C}_\OLId{latin-ordinary}{n}$، بحيث
تكون~$\Pi_\OLMathNumeral{0} \Sequent \Lambda_\OLMathNumeral{0}$ هي المتتالية النهائية
$\OLId{greek-variable}{Gamma} \Sequent \OLId{greek-variable}{Delta}$، وتكون~$\Pi_{\OLId{latin-ordinary}{n}+\OLMathNumeral{1}} \Sequent \Lambda_{\OLId{latin-ordinary}{n}+\OLMathNumeral{1}}$ مقدمةً
لـ~$\Pi_\OLId{latin-ordinary}{n} \Sequent \Lambda_\OLId{latin-ordinary}{n}$ وُلّدت في المرحلة~$\OLId{latin-ordinary}{n}+\OLMathNumeral{1}$ ولم تجد الخوارزمية
لها~!!a{proof}. ولا بد من مقدمة كهذه لكل~$\OLId{latin-ordinary}{n}$ ليس أخيرًا في الفرع، وإلا
لاهتدت الخوارزمية إلى~!!a{proof}. وإذا كان الفرع منتهيًا كانت
المتتالية الأخيرة~$\Pi_\OLId{latin-ordinary}{n} \Sequent \Lambda_\OLId{latin-ordinary}{n}$ موضع الإجهاض. ونجعل~$\OLId{latin-looped}{C}_\OLId{latin-ordinary}{n}$
مجموعة~!!{constant}s المرتبطة بالمتتالية~$\Pi_\OLId{latin-ordinary}{n} \Sequent \Lambda_\OLId{latin-ordinary}{n}$.

فإذا أجهضت الخوارزمية عند متتالية عليا ليس فيها إلا~!!{formula}s ذرية،
كان الفرع المنتهي بها فرع فشل. وإن لم تتوقف، توالت منها أشجار أكبر من
سابقاتها. ونجمع تصورها في شجرة نامية لامتناهية، هي ما كان يحصل لو نفذت
الخوارزمية عددًا لامتناهيًا من الخطوات. وهذه الشجرة، وإن كانت لامتناهية،
متناهية التشعب؛ فلكل عقدة ابن أو ابنان، هما المتتاليتان فوقها مباشرة.
وتقضي مبرهنة كونيغ بأن لكل شجرة لامتناهية متناهية التشعب فرعًا لامتناهيًا؛
وهذا الفرع هو فرع الفشل حين يدوم البحث أبدًا.

إذا كانت متتالية~$\Pi_\OLId{latin-ordinary}{n} \Sequent \Lambda_\OLId{latin-ordinary}{n}$ و$\OLId{latin-looped}{C}_\OLId{latin-ordinary}{n}$ فرع فشل، فلتكن
$\Theta=\bigcup_\OLId{latin-ordinary}{n}\Pi_\OLId{latin-ordinary}{n}$ و$\Xi=\bigcup_\OLId{latin-ordinary}{n}\Lambda_\OLId{latin-ordinary}{n}$ (بعد حذف المؤشرات
والوقائع المتكررة لـ~!!{formula}s)، ولتكن~$\OLId{latin-looped}{C}=\bigcup_\OLId{latin-ordinary}{n} \OLId{latin-looped}{C}_\OLId{latin-ordinary}{n}$.

وبهذا تهيأت مادة تعريف~$\Struct{M}$.
\begin{enumerate}
\item المجال~$\Domain{\OLId{latin-looped}{M}}$ هو مجموعة الحدود التي يمكن تكوينها من~$\OLId{latin-looped}{C}$ ومن
!!{function}s الواقعة في~$\OLId{greek-variable}{Gamma} \Sequent \OLId{greek-variable}{Delta}$ إن وُجدت، بغير أي
!!{variable}s.
\item إذا كان~$\OLId{latin-ordinary}{c}\in \OLId{latin-looped}{C}$، فإن~$\Assign{\OLId{latin-ordinary}{c}}{\OLId{latin-looped}{M}}=\OLId{latin-ordinary}{c}$.
\item إذا كان~$\OLId{latin-ordinary}{f}$ !!{function} ذا~$\OLId{latin-ordinary}{m}$ مواضع في~$\OLId{greek-variable}{Gamma} \Sequent \OLId{greek-variable}{Delta}$،
وكان~$\OLId{latin-ordinary}{t}_\OLMathNumeral{1},\dots,\OLId{latin-ordinary}{t}_\OLId{latin-ordinary}{m}\in\Domain{\OLId{latin-looped}{M}}$، فإن
$\Assign{\OLId{latin-ordinary}{f}}{\OLId{latin-looped}{M}}(\OLId{latin-ordinary}{t}_\OLMathNumeral{1},\dots,\OLId{latin-ordinary}{t}_\OLId{latin-ordinary}{m})=\OLId{latin-ordinary}{f}(\OLId{latin-ordinary}{t}_\OLMathNumeral{1},\dots,\OLId{latin-ordinary}{t}_\OLId{latin-ordinary}{m})$.
\item إذا كان~$\OLId{latin-looped}{R}$ !!{predicate} ذا~$\OLId{latin-ordinary}{m}$ مواضع في
$\OLId{greek-variable}{Gamma} \Sequent \OLId{greek-variable}{Delta}$، فإن
$\Assign{\OLId{latin-looped}{R}}{\OLId{latin-looped}{M}}=\Setabs{\tuple{\OLId{latin-ordinary}{t}_\OLMathNumeral{1},\dots,\OLId{latin-ordinary}{t}_\OLId{latin-ordinary}{m}}\in\Domain{\OLId{latin-looped}{M}}^\OLId{latin-ordinary}{m}}
{\OLId{latin-looped}{R}(\OLId{latin-ordinary}{t}_\OLMathNumeral{1},\dots,\OLId{latin-ordinary}{t}_\OLId{latin-ordinary}{m})\in\Theta}$.
\end{enumerate}
فـ~$\Struct{M}$ \emph{نموذج حدود}: مجاله مجموعة حدود، وتفسر فيه
!!{constant}s و!!{function}s بما يضمن~$\Value{\OLId{latin-ordinary}{t}}{\OLId{latin-looped}{M}}=\OLId{latin-ordinary}{t}$. ثم تستعمل
!!{formula}s الذرية في~$\Theta$ لتعريف~$\Assign{\OLId{latin-looped}{R}}{\OLId{latin-looped}{M}}$ على وجه يضمن أن
$\Sat{\OLId{latin-looped}{M}}{\OLId{latin-looped}{R}(\OLId{latin-ordinary}{t}_\OLMathNumeral{1},\dots,\OLId{latin-ordinary}{t}_\OLId{latin-ordinary}{m})}$ إذا وفقط إذا كان
$\OLId{latin-looped}{R}(\OLId{latin-ordinary}{t}_\OLMathNumeral{1},\dots,\OLId{latin-ordinary}{t}_\OLId{latin-ordinary}{m})\in\Theta$.

\begin{lem}\ollabel{lem:truth}
لكل~$!A\in\Theta\cup\Xi$:
\begin{enumerate}
  \item إذا كان~$!A\in\Theta$، فإن~$\Sat{\OLId{latin-looped}{M}}{!A}$.
  \item إذا كان~$!A\in\Xi$، فإن~$\Sat/{\OLId{latin-looped}{M}}{!A}$.
\end{enumerate}
\end{lem}

\begin{proof}
  نسلك الاستقراء على~$\depth{!A}$.

افترض أولًا أن~$!A$ ذرية؛ أي إما~$!A\ident\lfalse$ وإما
$!A\ident \OLId{latin-looped}{R}(\OLId{latin-ordinary}{t}_\OLMathNumeral{1},\dots,\OLId{latin-ordinary}{t}_\OLId{latin-ordinary}{m})$.

لأن~$\Pi_\OLId{latin-ordinary}{n} \Sequent \Lambda_\OLId{latin-ordinary}{n}$ فرع فشل، لا يجوز أن تحتوي أي~$\Pi_\OLId{latin-ordinary}{n}$ على
$\lfalse$ (وإلا كانت~$\Pi_\OLId{latin-ordinary}{n} \Sequent \Lambda_\OLId{latin-ordinary}{n}$ مسلّمة). وبحسب تعريف
$\Struct{M}$، يكون~$\Sat{\OLId{latin-looped}{M}}{\OLId{latin-looped}{R}(\OLId{latin-ordinary}{t}_\OLMathNumeral{1},\dots,\OLId{latin-ordinary}{t}_\OLId{latin-ordinary}{m})}$ إذا وفقط إذا كان
$\OLId{latin-looped}{R}(\OLId{latin-ordinary}{t}_\OLMathNumeral{1},\dots,\OLId{latin-ordinary}{t}_\OLId{latin-ordinary}{m})\in\Theta$. وتنتج عن الأمرين معًا العبارة: إذا كان
$!A\in\Theta$، فإن~$\Sat{\OLId{latin-looped}{M}}{!A}$.

إذا كانت~$!A$ ذرية و$!A\in\Pi_\OLId{latin-ordinary}{n}$، فإن~$!A\in\Pi_{\OLId{latin-ordinary}{n}+\OLMathNumeral{1}}$؛ وإذا كان
$!A\in\Lambda_\OLId{latin-ordinary}{n}$، فإن~$!A\in\Lambda_{\OLId{latin-ordinary}{n}+\OLMathNumeral{1}}$. (وهذا ناتج من الطريقة التي
تولّد بها خوارزمية البحث المتتاليات اللاحقة.) لذلك، إذا كان
$!A\in\Theta\cap\Xi$، كان~$!A\in\Pi_\OLId{latin-ordinary}{n}\cap\Lambda_\OLId{latin-ordinary}{n}$ لبعض~$\OLId{latin-ordinary}{n}$. وهذا يعني
أن~$\Pi_\OLId{latin-ordinary}{n} \Sequent \Lambda_\OLId{latin-ordinary}{n}$ مسلّمة، مع أن فرع الفشل لا يجوز أن يحتوي
مسلّمات. إذن، بحسب الإنشاء، $!A\notin\Theta\cap\Xi$ لكل~$!A$ ذرية؛ ومن ثم
إذا كان~$!A\in\Xi$ كان~$!A\notin\Theta$. ويستلزم تعريف~$\Struct{M}$ عندئذ
أن~$\Sat/{\OLId{latin-looped}{M}}{!A}$.

وفي خطوة الاستقراء نفرض صحة (\OLClassicalDigits{1}) و(\OLClassicalDigits{2}) لكل~!!{formula}s ذات~!!{depth}
أدنى من~$!A$، ثم نثبت (\OLClassicalDigits{1}) و(\OLClassicalDigits{2}) لـ~$!A$ بحسب الحالات.

لاحظ أولًا أنه إذا وقع~$!A^\OLId{latin-ordinary}{i}$ غير ذري في
$\Pi_\OLId{latin-ordinary}{n} \Sequent \Lambda_\OLId{latin-ordinary}{n}$، فإنه يقع كذلك في
$\Pi_{\OLId{latin-ordinary}{n}+\OLMathNumeral{1}} \Sequent \Lambda_{\OLId{latin-ordinary}{n}+\OLMathNumeral{1}}$ ما لم يكن~$\OLId{latin-ordinary}{i}$ أصغر مؤشر بين وقائع
الصيغ غير الذرية في~$\Pi_\OLId{latin-ordinary}{n} \Sequent \Lambda_\OLId{latin-ordinary}{n}$. وكل مؤشر جديد تضيفه
الخوارزمية أكبر من جميع المؤشرات الموجودة، ولذلك نصل في النهاية إلى
$\Pi_\OLId{latin-ordinary}{n} \Sequent \Lambda_\OLId{latin-ordinary}{n}$ يقع فيها~$!A^\OLId{latin-ordinary}{i}$ ويكون~$\OLId{latin-ordinary}{i}$ أصغر مؤشر بين
وقائع الصيغ غير الذرية. وفي المرحلة~$\OLId{latin-ordinary}{n}+\OLMathNumeral{1}$ تختزل الخوارزمية~$!A^\OLId{latin-ordinary}{i}$.
وبعبارة أخرى، إذا كان~$!A\in\Theta$، فلبعض~$\OLId{latin-ordinary}{n}$ يكون
$\Pi_\OLId{latin-ordinary}{n}=\Pi_\OLId{latin-ordinary}{n}',!A^\OLId{latin-ordinary}{i}$ ويكون~$\OLId{latin-ordinary}{i}$ أصغر مؤشر غير ذري. وإذا كان~$!A\in\Xi$،
فلبعض~$\OLId{latin-ordinary}{n}$ يكون~$\Lambda_\OLId{latin-ordinary}{n}=\Lambda_\OLId{latin-ordinary}{n}',!A^\OLId{latin-ordinary}{i}$ ويكون~$\OLId{latin-ordinary}{i}$ أصغر مؤشر غير ذري.

\begin{enumerate}
  \item إذا كان~$!A\in\Theta$ و$!A\ident !B\land !C$، فلبعض~$\OLId{latin-ordinary}{n}$ يكون
  $\Pi_\OLId{latin-ordinary}{n}=\Pi_\OLId{latin-ordinary}{n}',(!B\land !C)^\OLId{latin-ordinary}{i}$. وتكون
  $\Pi_{\OLId{latin-ordinary}{n}+\OLMathNumeral{1}}\Sequent\Lambda_{\OLId{latin-ordinary}{n}+\OLMathNumeral{1}}$ المقدمة المقابلة لـ~\LeftR{\land}؛ أي
  $\Pi_{\OLId{latin-ordinary}{n}+\OLMathNumeral{1}}=\Pi_\OLId{latin-ordinary}{n}',!B^\OLId{latin-ordinary}{k},!C^{\OLId{latin-ordinary}{k}+\OLMathNumeral{1}}$. ومن ثم~$!B,!C\in\Theta$. وبفرض
  الاستقراء~$\Sat{\OLId{latin-looped}{M}}{!B}$ و$\Sat{\OLId{latin-looped}{M}}{!C}$. وينتج من تعريف~$\Sat{\OLId{latin-looped}{M}}{}$ أن
  $\Sat{\OLId{latin-looped}{M}}{!B\land !C}$.

  \item إذا كان~$!A\in\Xi$ و$!A\ident !B\land !C$، فلبعض~$\OLId{latin-ordinary}{n}$ يكون
  $\Lambda_\OLId{latin-ordinary}{n}=\Lambda_\OLId{latin-ordinary}{n}',!B\land !C$. وتكون
  $\Pi_{\OLId{latin-ordinary}{n}+\OLMathNumeral{1}}\Sequent\Lambda_{\OLId{latin-ordinary}{n}+\OLMathNumeral{1}}$ إحدى مقدمتي~\RightR{\land} اللتين
  نتيجتهما~$\Pi_\OLId{latin-ordinary}{n}\Sequent\Lambda_\OLId{latin-ordinary}{n}',!B\land !C$؛ أي إما
  $\Lambda_{\OLId{latin-ordinary}{n}+\OLMathNumeral{1}}=\Lambda_\OLId{latin-ordinary}{n}',!B^\OLId{latin-ordinary}{k}$ وإما
  $\Lambda_{\OLId{latin-ordinary}{n}+\OLMathNumeral{1}}=\Lambda_\OLId{latin-ordinary}{n}',!C^\OLId{latin-ordinary}{k}$. ومن ثم إما~$!B\in\Xi$ وإما
  $!C\in\Xi$. وبفرض الاستقراء، إما~$\Sat/{\OLId{latin-looped}{M}}{!B}$ وإما
  $\Sat/{\OLId{latin-looped}{M}}{!C}$. وينتج من تعريف~$\Sat{\OLId{latin-looped}{M}}{}$ أن
  $\Sat/{\OLId{latin-looped}{M}}{!B\land !C}$.
  
  \item حالات~$!A\ident\lnot !B$ و$!A\ident !B\lor !C$ و
  $!A\ident !B\lif !C$ مماثلة، وتُترك تمارين.

  \item إذا كان~$!A\in\Theta$ و$!A\ident\lexists[\OLId{latin-ordinary}{x}][!B(\OLId{latin-ordinary}{x})]$، فلبعض~$\OLId{latin-ordinary}{n}$
  يكون~$\Pi_\OLId{latin-ordinary}{n}=\Pi_\OLId{latin-ordinary}{n}',\lexists[\OLId{latin-ordinary}{x}][!B(\OLId{latin-ordinary}{x})]^\OLId{latin-ordinary}{i}$. وتكون
  $\Pi_{\OLId{latin-ordinary}{n}+\OLMathNumeral{1}}\Sequent\Lambda_{\OLId{latin-ordinary}{n}+\OLMathNumeral{1}}$ المقدمة المقابلة لـ
  \LeftR{\lexists}؛ أي~$\Pi_{\OLId{latin-ordinary}{n}+\OLMathNumeral{1}}=\Pi_\OLId{latin-ordinary}{n}',!B(\OLId{latin-ordinary}{c})^\OLId{latin-ordinary}{k}$ لثابت~$\OLId{latin-ordinary}{c}$. ومن ثم
  $!B(\OLId{latin-ordinary}{c})\in\Theta$ و$\OLId{latin-ordinary}{c}\in \OLId{latin-looped}{C}$. وبفرض الاستقراء~$\Sat{\OLId{latin-looped}{M}}{!B(\OLId{latin-ordinary}{c})}$، ومن تعريف
  $\Sat{\OLId{latin-looped}{M}}{}$ ينتج لنا~$\Sat{\OLId{latin-looped}{M}}{\lexists[\OLId{latin-ordinary}{x}][!B(\OLId{latin-ordinary}{x})]}$.

  \item إذا كان~$!A\in\Xi$ و$!A\ident\lexists[\OLId{latin-ordinary}{x}][!B(\OLId{latin-ordinary}{x})]$، فلبعض~$\OLId{latin-ordinary}{n}$ يكون
  $\Lambda_\OLId{latin-ordinary}{n}=\Lambda_\OLId{latin-ordinary}{n}',\lexists[\OLId{latin-ordinary}{x}][!B(\OLId{latin-ordinary}{x})]^\OLId{latin-ordinary}{i}$. ولأن
  $\lexists[\OLId{latin-ordinary}{x}][!B(\OLId{latin-ordinary}{x})]$ يبقى في يمين المتتالية المقدمة بمؤشر جديد كلما
  اختُزل، فإنه يُختزل عددًا لامتناهيًا من المرات في يمين فرع الفشل متى وقع
  هناك. وكل حد~$\OLId{latin-ordinary}{t}\in\Domain{\OLId{latin-looped}{M}}$ يصير في النهاية ``أول حد لا يحتوي إلا
  ثوابت في~$\OLId{latin-looped}{C}_\OLId{latin-ordinary}{n}$ ولم يُستعمل من قبل في اختزال~$\lexists[\OLId{latin-ordinary}{x}][!B(\OLId{latin-ordinary}{x})]$ على
  اليمين'' على فرع الفشل. ولذلك، لكل~$\OLId{latin-ordinary}{t}\in\Domain{\OLId{latin-looped}{M}}$، يكون
  $!B(\OLId{latin-ordinary}{t})\in\Lambda_\OLId{latin-ordinary}{n}$ لبعض~$\OLId{latin-ordinary}{n}$، ومن ثم~$!B(\OLId{latin-ordinary}{t})\in\Xi$. وبفرض الاستقراء
  $\Sat/{\OLId{latin-looped}{M}}{!B(\OLId{latin-ordinary}{t})}$ لكل~$\OLId{latin-ordinary}{t}\in\Domain{\OLId{latin-looped}{M}}$. وينتج من تعريف~$\Sat{\OLId{latin-looped}{M}}{}$ أن
  $\Sat/{\OLId{latin-looped}{M}}{\lexists[\OLId{latin-ordinary}{x}][!B(\OLId{latin-ordinary}{x})]}$.

  \item حالات~$!A\ident\lforall[\OLId{latin-ordinary}{x}][!B(\OLId{latin-ordinary}{x})]$ مماثلة، وتُترك تمارين.
\end{enumerate}
\end{proof}

\begin{cor}\ollabel{cor:complete}
خوارزمية البحث عن البرهان في~\Log{G3c} كاملة: فإن أجهضت أو لم تتوقف،
كانت المتتالية الابتدائية~$\OLId{greek-variable}{Gamma} \Sequent \OLId{greek-variable}{Delta}$ غير صالحة، فلا
!!{proof} لها.
\end{cor}

\begin{proof}
  إذا أجهضت الخوارزمية أو لم تتوقف، وجد فرع فشل نعرف منه
  $\Struct{M}$. ومن~\olref{lem:truth} يكون~$\Sat{\OLId{latin-looped}{M}}{!A}$ لكل
  $!A\in\OLId{greek-variable}{Gamma}$ و$\Sat/{\OLId{latin-looped}{M}}{!A}$ لكل~$!A\in\OLId{greek-variable}{Delta}$. (تذكّر أن
  $\Pi_\OLMathNumeral{0}=\OLId{greek-variable}{Gamma}$ و$\Lambda_\OLMathNumeral{0}=\OLId{greek-variable}{Delta}$، ومن ثم~$\OLId{greek-variable}{Gamma}\subseteq\Theta$ و
  $\OLId{greek-variable}{Delta}\subseteq\Xi$.) وعلى هذا تكون~$\OLId{greek-variable}{Gamma} \Sequent \OLId{greek-variable}{Delta}$ غير صالحة.
  وبسلامة~\Log{G3c}، لا يوجد لها~!!{proof}.
\end{proof}

\begin{cor}
النسق~\Log{G3c} كامل: إذا كانت~$\OLId{greek-variable}{Gamma} \Sequent \OLId{greek-variable}{Delta}$ صالحة، فإن
$\Log{G3c}\Proves\OLId{greek-variable}{Gamma} \Sequent \OLId{greek-variable}{Delta}$.
\end{cor}

\begin{proof}
  نثبت المقابل النقيض. لنفرض أن
  $\Log{G3c}\Proves/\OLId{greek-variable}{Gamma} \Sequent \OLId{greek-variable}{Delta}$. لا بد أن تُجهض خوارزمية
  البحث عن البرهان أو ألا تتوقف، إذ لا يوجد~!!{proof} يمكن العثور عليه.
  وبموجب~\olref{cor:complete}، تكون~$\OLId{greek-variable}{Gamma} \Sequent \OLId{greek-variable}{Delta}$ غير صالحة.
\end{proof}

\end{document}
